\chapter{Problem Analysis}\label{ch:problem}
% Owner: M1 (Tushar B). Draft 1, 3 Nov 2026; sources: requirements (Project Overview,
% Problem Scenario) and blueprint §1.

\section{Problem}

A university campus depends on a set of shared IT services: Wi-Fi and the wired LAN, campus
email, the learning management system (LMS), laboratory PCs, printers and the Fee Payment
Portal. Thousands of students, faculty and staff use them every day, and a small IT helpdesk
keeps them running. When one of these services fails, the people affected need a reliable way
to report it, and the helpdesk needs to identify, prioritise, assign, investigate, resolve and
close the incident quickly.

Today that work is largely manual: problems arrive by phone, email or in person, are logged
inconsistently, and are prioritised and routed by judgement. The requirements list what this
leads to:

\begin{itemize}
  \item delayed responses, because nobody is told promptly that a ticket is waiting;
  \item incorrect prioritisation, because whoever shouts loudest is served first;
  \item inefficient technician allocation, because skills and current workload are not visible;
  \item SLA violations that are noticed only after the fact, if at all;
  \item poor visibility into service health;
  \item duplicate reports of the same outage, each handled separately;
  \item incomplete incident history, so nobody can say who did what and when;
  \item difficulty identifying recurring problems.
\end{itemize}

The requirements' central scenario makes this concrete. In our campus framing it reads:
\emph{the Fee Payment Portal is down and students cannot complete their fee transactions}.
Many students report it within minutes, the deadline for fee payment may be near, and the
outage may in fact be caused by a service the portal depends on, such as the database. The
system must turn that first report into a prioritised, assigned and monitored incident, keep
everyone informed, and record every step until the fix is verified and the incident closed.

\section{Stakeholders}

\noindent\begin{tabularx}{\textwidth}{@{}>{\raggedright\arraybackslash}p{4.2cm}>{\raggedright\arraybackslash}X@{}}
\toprule
\textbf{Stakeholder} & \textbf{Interest in the system} \\
\midrule
Employees (students, faculty, staff) & Report problems easily, see what happens to their report,
  answer questions from the helpdesk, and confirm that the fix works. \\
Technicians & See a clear queue of incidents that match their skills, without being overloaded;
  record investigation notes and resolutions. \\
Team Leads & Watch their team's queue and workload, take escalated incidents and reassign work. \\
Administrators & Configure users, technicians, services, categories, SLA policies and the
  assignment strategy; monitor incidents, SLA violations and the audit log. \\
IT Manager & The final escalation level; reads reports on SLA compliance, workload and service
  health to plan staffing and improvements. \\
Campus administration (CCIT, maintenance division) & Reliable services, measurable support
  quality, and an auditable record of incidents. \\
\bottomrule
\end{tabularx}

\section{Existing challenges}

Each challenge from the requirements maps to a cause and to the part of I-ITSM that addresses
it.

\noindent\begin{tabularx}{\textwidth}{@{}>{\raggedright\arraybackslash}p{3cm}>{\raggedright\arraybackslash}p{4cm}>{\raggedright\arraybackslash}X@{}}
\toprule
\textbf{Challenge} & \textbf{Root cause} & \textbf{Our response} \\
\midrule
Delayed response & No alerting; no response deadline & Response SLA per priority and
  service tier, notifications on creation and assignment, automatic escalation when nobody
  acknowledges. \\
Incorrect prioritisation & Reporters set priority; no rule & Priority computed from impact,
  urgency and service criticality; reporters answer plain-language questions; overrides
  need a reason and are audited. \\
Inefficient allocation & Skills and load are invisible & Assignment strategies that weigh
  skills, priority-weighted workload, experience and past fixes on the service. \\
SLA violations & Deadlines are not tracked & An SLA clock per incident (paused while waiting on
  the reporter), a background monitor, warnings before the deadline and recorded violations. \\
Poor visibility & No shared view of service state & Live dashboards, a live service health map, and
  a dependency graph that names the likely root cause of related incidents. \\
Duplicate reports & Each report is handled on its own & Hybrid duplicate detection while the
  reporter types, and linking of confirmed duplicates to the original. \\
Incomplete history & Actions are not recorded & An append-only, hash-chained audit trail of
  every state change, login and configuration change. \\
Recurring problems & History is not searchable & Nine-criteria search, analytics and trends,
  suggested resolutions from similar past incidents, and correlation through service
  dependencies. \\
\bottomrule
\end{tabularx}

\section{Proposed solution}

I-ITSM is a full-stack web application that manages the whole incident lifecycle:
\textbf{Report $\rightarrow$ Classify $\rightarrow$ Prioritise $\rightarrow$ Assign
$\rightarrow$ Investigate $\rightarrow$ Resolve $\rightarrow$ Verify $\rightarrow$ Close}.
It has three visibly different portals: Employee, Technician and Administrator.

\begin{itemize}
  \item \textbf{Lifecycle.} Ten states and exactly twenty legal transitions, implemented with
    the State pattern, so illegal moves such as NEW $\rightarrow$ CLOSED are impossible. Hold,
    reopen, cancel and a 72-hour auto-close are part of the lifecycle.
  \item \textbf{Core engines.} Priority ($S = I \times U \times C$, mapped to P1--P4), SLA targets
    by service tier and priority, three interchangeable assignment strategies, notifications
    through several channels, and a Chain-of-Responsibility escalation from technician to
    Team Lead, Administrator and IT Manager.
  \item \textbf{Intelligent features.} AI-assisted classification with a rule-based fallback,
    suggested resolutions from similar resolved incidents, hybrid duplicate detection,
    dependency-based correlation, predictive SLA-risk estimation and smart technician
    recommendation with reasons.
  \item \textbf{Architecture.} A layered modular monolith whose framework-free domain model is
    the centre of the design; the API, persistence and external services depend on it, never
    the other way round (Chapter~\ref{ch:architecture}).
\end{itemize}

In the Fee Payment Portal scenario, the employee's report creates an incident; the system
suggests the Web Application category, identifies the Gold-tier portal (criticality 3), computes
$3 \times 3 \times 3 = 27$, which is P1, applies the 5-minute response and 30-minute resolution
targets, recommends and assigns the best technician, notifies the reporter and the technician,
watches the deadlines, escalates if nobody acknowledges, links the other students' duplicate
reports, points to a likely root cause in a dependency, records every step, and finally counts
the incident in the reports once the reporter has verified the fix.

\section{Assumptions}

\begin{enumerate}
  \item One organisation (the campus), with incidents written in English.
  \item SLA clocks run 24$\times$7; business-hours calendars are out of scope.
  \item The historical incidents used for analytics, prediction and suggested resolutions are
    realistic synthetic seed data (420 incidents over four months).
  \item Reporters do not set priority directly; they answer two plain-language questions, and the
    service's criticality comes from its configuration.
  \item Silver and Bronze SLA targets are twice and four times the Gold targets given in the
    requirements.
  \item External services (Gemini, email, Telegram, Google Sign-In) can fail; the system then
    falls back to deterministic behaviour and never blocks incident creation.
  \item The API runs as a single instance, so the in-process SLA monitor needs no coordination;
    Chapter~\ref{ch:technology} explains how this changes at scale.
\end{enumerate}
